\chapter{Discussion and Conclusion}
\label{ch:discussion}

\section{Discussion}
\label{sec:discussion}

Overall the results show that users are accepting of the passwordless system as they generally find it usable, secure, and preferred to a password-based system especially when sharing many files or sharing files with individuals external to a controlled environment. Of the hypotheses stated in Section~\ref{sec:hypothesis}, we find that users perceived greater security, comparable usability, and better deployability in the passwordless system compared to the password-based approach. Both systems scored within the ``Acceptable'' range of usability on the SUS (72.0 vs.\ 73.3, $p=.944$), and the passwordless system was rated significantly more useful on the Van der Laan acceptance scale ($p=.014$), with a similar trend toward higher satisfaction ($p=.078$). That a novel prototype achieved usability parity with a mature commercial tool on first use is itself notable, as refinements to the user interface and increased familiarity would likely improve the passwordless system's scores further. Qualitatively, users identified unique benefits of the passwordless system more often than they shared a preference for the password-based system. However, some users did note a preference for the password-based system, typically citing its simplicity, familiarity, and ease of use on the receiving side. Nearly all participants (20 of 21, see Table~\ref{tab:willingness_to_adopt}) stated they would use the passwordless system for their own documents if it were available today, and when asked which method they would prefer for sharing multiple documents at scale, 19 of 21 chose the passwordless approach. Participants frequently saw the passwordless method as more trustworthy and secure than traditional password-based approaches, citing frustration with password sharing and trust in biometric authentication as key factors.

A key finding from the qualitative results is that participants widely disliked and distrusted passwords while expressing strong preference for biometric authentication. Participants frequently noted frustration with creating and sharing passwords in the password-based process, while the biometric interaction in the passwordless process was met with immediate comfort: participants would tap their fingerprint without hesitation, often before fully reading the on-screen instructions.

Users intuitively associated passkey uniqueness with biometric uniqueness. Although the passkey is not derived from the user's biometric information, both share a uniqueness property that users naturally grasp. Users know ``I am the only one with my fingerprint and face'' but ``anyone could guess my password''. As one participant said, ``my face feels more secure than any password'' (P24). The biometric interaction in passkey creation thus serves as a tangible proxy for security that participants can understand without needing to understand the underlying cryptography. This mental model---``my fingerprint is unique, therefore my encryption is unique''---may be technically imprecise, but it maps onto the actual security properties closely enough to support accurate trust judgments. While the concept of a passkey is abstract and unfamiliar to most users, the biometric interaction is familiar, simple, and trusted. Several participants (10 of 21) mentioned a lack of understanding about how passkeys work technically with some citing it as a reason for hesitation in trusting the passwordless system. However, most participants, including many who felt the passkey process was abstracted, reiterated that the passwordless method, anchored on their biometric action, felt more secure.

Several factors may contribute to the perceived security of the passwordless method. Some participants' trust appeared rooted in the biometric interaction itself: the intuitive uniqueness of a fingerprint or face. Others grounded their trust in the platform provider's brand (P17: ``If Adobe had a passwordless one, I trust Adobe''; P29: ``if Google does it, it's fine''). Still others may associate the novelty or sophistication of the interaction with stronger security. Disentangling these factors is difficult with the current data, but the consistency of the ``feels more secure'' sentiment across participants with different reasoning suggests that the biometric interaction serves as a unifying trust signal regardless of the underlying rationale.

These findings align with recent work by Reittinger et al.~\cite{reittinger_passkey_2025}, who found a significant increase in passkey adoption when applying digital nudges that emphasized biometric framing. For example, they described passkeys as ``secure keys that use your fingerprint, face, or screen lock to simplify access''~\cite{reittinger_passkey_2025}. Our participants' behavior confirms this approach: the biometric interaction served as an implicit nudge, making the security action feel familiar and trustworthy without requiring users to understand the underlying passkey mechanism. These findings suggest that passwordless adoption is most likely to succeed when the system emphasizes the biometric interaction and focuses on simplicity, ease of use, and uniqueness rather than the technical specifics of FIDO2 and WebAuthn.

The biometric interaction also connects to the self-efficacy literature discussed in Section~\ref{sec:quest-passwords}. Biometric interaction provides the high coping appraisal identified as the strongest predictor of security behavior adoption: users feel capable of performing the security action---a fingerprint tap is something they already do daily---and believe it effectively protects their data because the uniqueness of their biometric maps intuitively onto the uniqueness of the encryption key. This combination of self-efficacy and response efficacy---``I can do this, and it works''---stands in contrast to the password-based method, where participants expressed low confidence in their ability to create strong passwords, remember them, and share them securely. Security tools that make users feel confident in performing the secure action are more likely to drive adoption than tools that emphasize the severity of threats. While the series of MFT breaches described in Chapter~\ref{ch:background} provides strong motivation for change in document protection methods, users are more likely to respond to improved tools than to threat awareness alone. Insecure behavior---in this case using simple, repeated passwords or sharing the password through an insecure method---is mostly driven by the security process being too effortful or complex. Users are often rational when balancing risk against effort and understand that worst-case harm is not the same as actual harm~\cite{herley_so_2009}. Human effort and attention is a precious resource primarily dedicated to productivity, so security tools need to fit into work processes without adding extra friction~\cite{enisa_cybersecurity_2019}. In the case of document encryption and sharing, the incumbent password-based encryption method already presents exactly this kind of friction, and participants in this study were eager to adopt methods that reduce the time, effort, and attention needed to share a file securely.

Usability challenges remain with the passwordless system presented in this work, which is in a proof-of-concept state and was not tested across every browser and device operating system due to the BYOD nature of the study. However, even in its early iteration, the passwordless system's usability friction was comparable to that of the established password-based method: three of 21 participants encountered problems creating or using passkeys, while three of 21 participants encountered problems sharing the password. The types of friction differ: passkey issues were platform-specific and likely to improve as WebAuthn support matures, while password sharing challenges are inherent to the password-based model itself. This parity of usability issues observed in real participant behavior aligns with the parity of usability ratings at the SUS ``Acceptable'' level. The significant correlation between technology affinity and passwordless usability ($\rho=+.578$, $p=.006$) suggests that this parity was achieved with a technically literate sample; broader adoption among less technical populations will likely require additional onboarding support and interface simplification. Improving the user experience of the encryption flow and allowing FIDO2 protocols to gain more familiarity and compatibility offer room for the passwordless method to surpass password-based encryption in the near future. One notable tension emerged: P7 observed that the passwordless process felt ``almost too easy to be secure,'' suggesting that high ease of use can paradoxically undermine confidence when perceived effort does not match users' mental model of what security should feel like. Providing visible confirmation that the encryption succeeded---such as a clear file state change or explicit success message---may help bridge this gap between ease and trust.

As discussed in Chapter~\ref{ch:evaluation}, deployability concerns are a common reason novel security architectures fail to gain traction and widespread use. While deployability was examined as part of this study's evaluation framework, this work is primarily an early design and prototyping effort aimed at discovering whether passwordless document encryption is feasible and acceptable from a user perspective.

The adoption of SSL/TLS offers an instructive parallel when considering how novel security mechanisms can gain widespread adoption. Reflecting on the success of SSL/TLS adoption which is used today to encrypt web traffic by default, Garfinkel et al.\ identify three theories of disuse and solutions that made user adoption easier. Their theories of disuse: (1) that users lack the software, (2) that it is too difficult to use, and (3) that users simply do not want to use it, are addressed through (1) distribution with operating systems, (2) making it automatic, and (3) user education~\cite{garfinkel_email_2005}. These theories of disuse informed the design requirements of the passwordless document encryption system: it runs in a standard web browser requiring no software installation, automates key management through WebAuthn PRF rather than requiring manual password creation, and provides a familiar file-based workflow. While these design choices address the theories of disuse at the prototype level, whether they hold at scale in production environments remains an open question explored further in Section~\ref{sec:future-work}.

The three research questions which guided this work build from problem identification to system design and solution validation. The comparative framework (R1) identified the limitations of existing document encryption methods and derived design requirements for an improved system. Key limitations included a perimeter defense model, the cognitive burden of creating and sharing secrets, vulnerability to brute-force attacks, inability to recover from a forgotten password, and the requirement for recipients to install specific software. The system design (R2) addressed those requirements through browser-based age encryption with WebAuthn PRF key derivation, providing strong modern encryption with cryptographic agility, keeping all plaintext data and encryption keys local to the browser, and requiring no passwords or software installation by either sender or recipient. These properties directly address the architectural vulnerabilities exposed by the series of MFT breaches described in Chapter~\ref{ch:background}, where compromised server sessions exposed every document on the platform. The user study (R3) provides initial empirical evidence that this approach is viable with participants finding the passwordless system usable and broadly preferring it to password-based encryption. While this work represents an early design and prototyping effort, it demonstrates that passwordless document encryption is both technically feasible and acceptable to users.

\section{Limitations}
\label{sec:limitations}

The study has several limitations primarily related to sample size, demographic diversity, and the experimental system design.

The participant pool (N=21) was relatively small and demographically homogeneous, skewing heavily male (76.2\%), relatively young ($M=29.2$, $SD=11.7$), and technically educated, with 71.4\% possessing a formal Computer Science or Cybersecurity background. While this provides valuable initial insights into usability among early adopters and tech-literate users, it limits the generalizability of the findings to the broader, non-technical public. Several participants explicitly noted that users without their technical background might not feel as confident with the passwordless system, underscoring this constraint.

The within-subjects design, where each participant evaluated both systems, introduces the possibility of order and contrast effects. This is a well-known limitation of within-subjects user studies where participants' experience with the first system influences their evaluation of the second. To mitigate this, the presentation order was counterbalanced: 11 participants evaluated the password-based Adobe system first, and 10 evaluated the prototype passwordless system first. This counterbalancing helps average out learning and order bias across the sample, but cannot eliminate contrast effects at the individual level, where participants may still anchor their assessment of one system against the other.

Additionally, with usability sessions lasting 30--35 minutes, learning fatigue may have influenced response quality in later portions of the study. To reduce fatigue, we intentionally placed the privacy scale (PS), affinity for technology use (ATI), demographics and other background questions in the pre-session survey to reduce the time answering survey questions during the usability session. Despite these measures, non-negligible learning fatigue may have influenced the quality of the results.

Finally, as with any prototype evaluation, the implemented system has not reached the maturity level of a production tool like Adobe Acrobat. The design builds on established standards---FIDO2, WebAuthn, and age encryption---and runs within a standard web browser, which partially mitigates the maturity gap compared to a fully novel prototype. However, there remains a likelihood that participants evaluated the prototype negatively based on its lack of polish rather than on the merits of the passwordless flow itself.

\section{Future Research Directions}
\label{sec:future-work}

Future work should address the methodological limitations identified in Section~\ref{sec:limitations}. A between-subjects design, in which participants evaluate only one encryption method, would eliminate contrast and order effects but would require a substantially larger sample size to achieve comparable statistical power. Recruitment should target a more demographically diverse population, including non-technical users, to assess whether the usability findings generalize beyond early adopters. Additionally, a future study could be conducted asynchronously, with participants evaluating the system on their own devices and reporting feedback through a survey instrument rather than an in-person interview, enabling broader geographic and demographic reach. Such a study could also target participants in regulated environments such as healthcare or finance, where document encryption is an operational requirement rather than an abstract task.

Beyond the user study, translating the prototype system into a production-ready implementation remains an open challenge. Future work could explore how passwordless document encryption might be distributed within operating systems and integrated with platform authenticators to perform per-file encryption automatically. As discussed in Section~\ref{sec:discussion}, automating the encryption process and embedding it within the operating system may remediate common reasons for user avoidance and---similar to SSL/TLS---help an improved security architecture become a default across many users and devices. Several directions to extend the technical architecture described in Chapter~\ref{ch:system-design} are described in the remainder of this section.

While typage was selected for this implementation because TypeScript interfaces directly with the WebAuthn and FIDO2 browser APIs, a WebAssembly (Wasm) build of the age encryption engine would offer additional security benefits through sandboxed execution, isolating the encryption process from the broader JavaScript runtime. Projects have demonstrated age encryption running as a Go build within Wasm~\cite{age_wasm}. A Wasm build of the current system would require bridging between the sandboxed encryption environment and the browser's credential APIs. A Wasm approach also opens the door to implementing the encryption process in a memory-safe language such as Rust or Go, which would reduce the attack surface associated with npm supply chain risks and JavaScript/TypeScript runtime vulnerabilities. This would further strengthen the system's security posture while still enabling the usability and deployability benefits of the system running in a web browser session.

Recent advances in quantum computing have compressed expert timelines for cryptographically-relevant quantum computers significantly, with Google targeting 2029~\cite{google_pqc_2026} and the Department of Defense aiming to fully implement post-quantum encryption by 2030~\cite{dod_pqc_2026}. This urgency is acute for file encryption systems, which are vulnerable to store-now-decrypt-later attacks that motivate a move to post-quantum encryption before attacks materialize. As of version 1.3.0, age natively supports hybrid post-quantum recipients using ML-KEM-768 combined with X25519. The prototype system exposes this as a selectable encryption mode (see Figure~\ref{fig:ui-settings}), but it cannot yet be combined with the WebAuthn PRF hardware-binding flow. Users must choose between passwordless hardware-bound encryption and post-quantum protection today. While the system's architecture was designed and built for cryptographic agility where the encryption method can be changed without lengthy system rebuilds, work remains to unify the WebAuthn PRF authentication ceremony with age's post-quantum recipients feature to ensure hardware-bound encryption that is also quantum-resistant. Specifically, the ECDSA signature in the WebAuthn authentication ceremony (usually ES256) should be replaced with a post-quantum secure implementation (although it is partially mitigated by the interactive, online nature of the ceremony) and the resulting age recipient processes should integrate the hybrid ML-KEM-768 and X25519 method to wrap the file based on the WebAuthn value returned.

More complete attribute-based access control through OpenTDF would enable policies to travel with the encrypted file itself, supporting fine-grained rules such as time-based expiration, role-based access, and attribute-level restrictions beyond the current recipient-list model. The OpenTDF-style manifest included in the current system provides a foundation that could be expanded to provide more robust access control supporting the use of the system in zero-trust networks. This future research direction could also address how recipient identities are associated with contact identifiers such as email addresses, and whether that association should be global or per-document (as discussed in Section~\ref{sec:request-share-flows}).

Additionally, integrating the system with cloud identity federation (e.g., AWS IAM/STS with OIDC) would enable the system to work with cloud storage services such as Box, Google Drive, or AWS S3. This approach would provide audit logging, short-term scoped credentials, and ABAC-governed access to cloud-hosted encrypted bundles. Access sharing would be administered through the cloud service's native access controls rather than standalone share links. The current system includes scaffolded share-link functionality to demonstrate the architectural pattern, but an improved implementation would integrate with established cloud storage providers.

\section{Conclusion}
\label{sec:conclusion}

Encrypted documents are used everyday in finance, legal, healthcare and other professions to protect sensitive information that must be shared with external partners. However, current document encryption methods do not adequately protect data confidentiality as demonstrated by the series of MFT breaches affecting millions of individuals described in Chapter~\ref{ch:background}. Password-based encryption remains the most widely available option, but it is vulnerable to brute-force attacks, unauthorized sharing, and the inherent challenge of securely communicating passwords to recipients.

This work presents an alternative approach. Through a structured comparative framework, we identified the usability, deployability, and security limitations of nine existing document encryption methods. From this analysis, we derived 15 design requirements and performed an iterative system design which resulted in a novel system built on FIDO2 passkeys, the WebAuthn PRF extension, and open-source age encryption. The system provides non-custodial end-to-end document protection without passwords removing the third-party vendor risk which caused the MFT breaches. We create a proof-of-concept implementation of our design as an open-source codebase and website. The system requires no software installation, derives encryption keys from hardware-bound authenticator secrets rather than user-chosen passwords, and performs all encryption locally in the browser.

We evaluated this system through a counterbalanced within-subjects user study (N=21) comparing the passwordless approach to traditional password-based encryption via Adobe Acrobat. We utilize the Systems Usability Scale (SUS) and Van der Laan Acceptance Scale as metrics for empirical analysis. The passwordless prototype achieved comparable usability on first use (SUS 72.0 vs.\ 73.3, $p=.944$) while being rated significantly more useful (VdL usefulness $p=.014$). Nineteen of 21 participants preferred the passwordless method for sharing documents at scale. Qualitative results show that users felt frustration with password sharing and felt a strong sense of security and trust in the biometric authentication used in the passwordless system. Together these factors contributed to a strong preference for using the passwordless approach.

To our knowledge, this is the first study to apply FIDO2 and WebAuthn PRF to end-to-end document encryption rather than user authentication, and one of the first human-subjects usability evaluations of the PRF extension in any context. This work offers several implications for practitioners and researchers.

For industry, the findings suggest that passwordless document encryption is both technically feasible and acceptable to users. Practitioners designing such systems should prioritize biometric-framed authentication flows, as users associate the uniqueness of fingerprint and facial recognition with security even without understanding the underlying cryptography. The password sharing burden---identified by a majority of participants as the weakest link in the traditional approach---should be eliminated. Systems that remove this step gain an immediate usability and security advantage.

For researchers, this study demonstrates that the WebAuthn PRF extension is a viable foundation for encryption applications beyond user authentication, opening a design space that has received little academic attention. The structured comparative framework presented in Chapter~\ref{ch:evaluation} provides a reusable tool for evaluating future document encryption methods. The significant correlation between ATI and passwordless usability ($p=.006$) suggests that future studies should recruit more demographically diverse samples and investigate onboarding interventions for less technical users.

Sensitive documents deserve the same level of cryptographic protection that has become standard for web traffic and instant messaging. This work provides initial evidence that passwordless, per-document encryption is a viable path toward that goal.
